\chapter{Jaundice in Adults}

\section*{Definition and Overview}
Jaundice is a yellow discolouration of the skin and the whites of the eyes caused by a build-up of bilirubin, a yellow pigment produced when red blood cells are broken down. It is a sign of an underlying problem, not a disease in itself. The causes fall into three groups: increased production of bilirubin, as in haemolysis; inability of the liver to process bilirubin, as in hepatitis, cirrhosis, or inherited conditions; and blockage of bile drainage, as in gallstones, bile duct tumours, and pancreatic cancer. Jaundice always warrants medical assessment, because some causes, such as bile duct obstruction, are treatable but can become life-threatening if delayed.

\section*{Epidemiology}
Jaundice in adults is less common than newborn jaundice and is always significant. Gallstones are a leading cause of obstructive jaundice, affecting up to one in ten Australians, with bile duct stones causing obstruction in a proportion. Viral hepatitis, including hepatitis A, B, and C, causes jaundice in some infections, and alcohol-related liver disease and cirrhosis account for many cases. Pancreatic cancer, which often presents with painless jaundice, affects around 3,000 Australians each year. Drug-induced liver injury from medicines and herbal products is an important and preventable cause.

\section*{Aetiology and Pathophysiology}
Bilirubin is produced from the breakdown of haemoglobin and is processed by the liver and excreted in bile. Jaundice appears when this system fails at any point. Pre-hepatic jaundice occurs when excessive red cell breakdown overwhelms the liver, as in haemolysis. Hepatic jaundice occurs when the liver cells cannot take up or process bilirubin, as in viral or alcoholic hepatitis, cirrhosis, drug-induced liver injury, and inherited conditions such as Gilbert's syndrome. Post-hepatic (obstructive) jaundice occurs when bile cannot drain, as in gallstones, bile duct strictures, and tumours of the bile duct or pancreas. In obstruction, the bile backs up, bilirubin rises, the stool becomes pale, and the urine darkens, and prolonged blockage damages the liver.

\section*{Clinical Features}
\begin{itemize}
    \item Yellowing of the skin and the whites of the eyes
    \item Dark urine, like tea or cola
    \item Pale, clay-coloured stools, suggesting obstruction
    \item Itching of the skin, common with bile obstruction
    \item Fatigue, malaise, and nausea
    \item Fever and abdominal pain, which may indicate infection or gallstones
    \item Weight loss, which may suggest cancer
    \item Easy bruising and bleeding in advanced liver disease
    \item Confusion or drowsiness in severe liver failure
\end{itemize}

\section*{Differential Diagnosis}
The cause of jaundice must be identified through investigation.
\begin{itemize}
    \item \textbf{Haemolysis:} increased red cell breakdown from blood disorders
    \item \textbf{Gallstones:} blockage of the bile duct, often with pain
    \item \textbf{Viral hepatitis:} hepatitis A, B, C, and E
    \item \textbf{Alcohol-related liver disease and cirrhosis:} chronic liver damage
    \item \textbf{Drug-induced liver injury:} medicines, including paracetamol, and herbal products
    \item \textbf{Bile duct and pancreatic cancer:} often painless jaundice with weight loss
    \item \textbf{Inherited conditions:} Gilbert's syndrome and other metabolic disorders
    \item \textbf{Other liver diseases:} autoimmune, fatty liver, and primary biliary cholangitis
\end{itemize}

\section*{When to Seek Medical Care}
Jaundice always requires medical assessment, and urgent assessment is needed for jaundice with abdominal pain, fever, confusion, bleeding, or signs of severe illness. Anyone with dark urine, pale stools, or unexplained itching should also be assessed. People at risk of hepatitis should be tested, and those taking medicines that affect the liver should report new symptoms.

\textbf{Call 000 (or local emergency services) for:}
\begin{itemize}
    \item Confusion, drowsiness, or altered consciousness, which may indicate liver failure
    \item Vomiting blood or black, tarry stools
    \item Severe abdominal pain with fever
    \item Difficulty breathing or collapse
    \item Bleeding that is hard to stop
\end{itemize}

\section*{Diagnosis and Investigations}
\begin{itemize}
    \item \textbf{Blood tests:} liver function tests, bilirubin, and clotting
    \item \textbf{Full blood count and blood film:} to assess for haemolysis
    \item \textbf{Hepatitis serology:} testing for hepatitis A, B, C, and E
    \item \textbf{Ultrasound:} first-line imaging of the liver, gallbladder, and bile ducts
    \item \textbf{CT or MRI:} detailed imaging for tumours and obstruction
    \item \textbf{ERCP or MRCP:} examination and treatment of the bile ducts
    \item \textbf{Liver biopsy:} in selected cases to diagnose liver disease
    \item \textbf{Medication review:} identification of drug-induced causes
\end{itemize}

\section*{Management}
\begin{itemize}
    \item \textbf{Treatment of the cause:} antibiotics and removal of gallstones, treatment of hepatitis, and stopping culprit medicines
    \item \textbf{Obstruction relief:} ERCP with stone removal or stenting of the bile duct
    \item \textbf{Surgery:} for tumours and complex obstruction
    \item \textbf{Management of liver disease:} alcohol cessation, treatment of cirrhosis, and management of complications
    \item \textbf{Supportive care:} fluids, nutrition, treatment of itching, and vitamin replacement
    \item \textbf{Liver transplant:} for advanced liver failure in suitable cases
    \item \textbf{Monitoring:} regular liver tests and imaging
    \item \textbf{Public health:} vaccination and testing for hepatitis
\end{itemize}

\section*{Complications}
\begin{itemize}
    \item Liver failure with confusion, bleeding, and kidney failure
    \item Infection of the bile ducts (cholangitis), which is life-threatening
    \item Acute pancreatitis from gallstones
    \item Cirrhosis and its complications, including liver cancer
    \item Malnutrition and vitamin deficiencies from bile obstruction
    \item Chronic itching and sleep disturbance
    \item Delayed diagnosis of cancer with worse outcomes
\end{itemize}

\section*{Prognosis and Outcomes}
The outlook depends entirely on the cause. Obstructive jaundice from gallstones resolves with stone removal, hepatitis A and E usually resolve, and stopping a culprit medicine reverses drug-induced injury. Chronic liver disease and cirrhosis can be slowed with treatment but are progressive, and pancreatic and bile duct cancer carry a serious prognosis that improves with early diagnosis. Because many causes are treatable, prompt investigation of jaundice is essential and improves outcomes.

\section*{Prevention and Self-Care}
\begin{itemize}
    \item Limit alcohol and maintain a healthy weight
    \item Be vaccinated against hepatitis A and B
    \item Practise safe injecting and safe sex to prevent hepatitis B and C
    \item Take medicines, including paracetamol, at recommended doses
    \item Avoid herbal products and supplements that can damage the liver
    \item Seek medical advice promptly for yellowing, dark urine, or pale stools
    \item Attend follow-up testing for liver disease
    \item Discuss pregnancy, medicines, and travel with your doctor if at risk
\end{itemize}

\section*{Sources and Further Reading}
The following reputable sources provide further information for healthcare students and patients seeking reliable, current advice about jaundice in adults.
\begin{itemize}
    \item Healthdirect Australia. \textit{Jaundice in adults.} https://www.healthdirect.gov.au/
    \item Hepatitis Australia. \textit{Hepatitis information and support.} https://www.hepatitisaustralia.com/
    \item Gastroenterological Society of Australia. \textit{Liver disease information.} https://www.gesa.org.au/
    \item Cancer Council Australia. \textit{Pancreatic cancer information.} https://www.cancer.org.au/
    \item NHS. \textit{Jaundice.} https://www.nhs.uk/conditions/jaundice/
    \item Mayo Clinic. \textit{Jaundice.} https://www.mayoclinic.org/
\end{itemize}
